%lezione 9 del 26 marzo 2026
\newpage
\section{Il propagatore di un'Hamiltoniana con potenziale}
Si consideri ora un sistema fisico a $d$ gradi di libertà con Hamiltoniana $H$ del tipo
$$H = \sum_{\alpha=1}^{d}\frac{p^2_\alpha}{2m} + V(q_\alpha)$$
Il propagatore tra $(q,t)$ e $(q',t')$ è
$$K(q',t',q,t) = \bra{q'}U(t',t)\ket{q} = \bra{q'}\exp\left\{\frac{-i}{\hbar}H(t'-t)\right\}\ket{q}$$
come per la particella libera, si tenta di aggiungere la relazione di completezza di $\ket{p}$
$$K(q',t',q,t) = \int d^dp\bra{q'}\ket{p}\bra{p}\exp\left\{\frac{-i}{\hbar}H(t'-t)\right\}\ket{q}$$
A questo punto però si presenta un problema non indifferente:
l'elemento di matrice $\bra{p}U\ket{q}$ non è facilmente determinabile come si è fatto in precedenza,
perché i termini non lineari dell'espansione dell'esponenziale non sono riordinabili
(essendo infiniti).

Bisogna dunque procedere diversamente, si divide l'intervallo di tempo $t'-t$ in $N$ intervalli di
lunghezza $\Delta t$ e si sfrutta la proprietà (\ref{eq:K_iter}) per scrivere il propagatore
$$K=\int \prod_{i=1}^{N-1} d^fq_i \prod_{j=0}^{N}\bra{q_{j+1}}U(t_{j+1},t_j)\ket{q_j}$$
con $t_0=t<t_1<...<t_N=t'$, $q_0=q$, $q_N=q'$ e $t_{j+1}-t_{j}=\Delta t$.\\
Se $\Delta t \to0$, si può linearizzare l'espressione di $U(t_{j+1},t_j)$ e procedere come
per la particella libera:
\begin{gather*}
\bra{q_{j+1}}U(t_{j+1},t_j)\ket{q_j}\sim \bra{q_{j+1}}\mathbb{I}-\frac{i}{\hbar}\left(\frac{\hat p^2}{2m}+V(\hat q)
\right)\Delta t\ket{q_j}=\\
=\int d^dp_j \bra{q_{j+1}}\ket{p_j}\bra{p_j}\mathbb{I}-\frac{i}{\hbar}\left(\frac{\hat p^2}{2m}+
V(\hat q)\right)\Delta t\ket{q_j}=\\
=\int d^dp_j \bra{q_{j+1}}\ket{p_j}\left[1-\frac{i}{\hbar}\Delta t\left(\frac{p_j^2}{2m}+V(q_j)
\right)\right]\bra{p_j}\ket{q_j}=\\
=\int \frac{d^dp_j}{(2\pi\hbar)^d}\exp\left\{\frac{i}{\hbar}\Delta t\left(\frac{p_j(q_{j+1}-q_j)}
{\Delta t} - H(p_j,q_j)\right)\right\}
\end{gather*}
dove si è ritornati alla forma esponenziale nell'ultimo passaggio.\\
Reinserendo questa espressione nel calcolo del propagatore, si ottiene infine
\begin{equation}
\boxed{K(q',t',q,t) = \int\prod_{j=1}^{N-1}d^dq_j\prod_{k=0}^{N-1}\frac{d^dp_k}{(2\pi\hbar)^d}
\exp\left\{\frac{i}{\hbar}\sum_{i=0}^{N-1}\Delta t \left(\frac{p_i(q_{i+1}-q_i)}{\Delta t}-H(p_i,q_i)
\right)\right\}}
\label{eq:K_discr}
\end{equation}
\textbf{Oss.} Questa espressione discretizzata per $K$ è un integrale su tutti i possibili punti
intermedi tra lo stato inizale e finale nello \textit{spazio delle fasi} del sistema.

Si denota ora $\dot q_i:=\frac{q_{i+1}-q_i}{\Delta t}$ e si considera il limite continuo
$n\to\infty$: l'argomento dell'esponenziale è una somma alla Riemann, che diventa un integrale sul
tempo, mentre l'integrale sui punti intermedi diventa l'integrale su tutti i possibili
\textit{cammini} nello spazio delle fasi e la sua misura di integrazione si denota con
$\mathcal{D}q\mathcal{D}p$
$$K(q',t',q,t) = \int\mathcal{D}q\mathcal{D}p\exp\left\{\frac{i}{\hbar}\int_{t}^{t'}dt(p\dot q - H)
\right\}$$
Questa espressione può essere ulteriormente semplificata per Hamiltoniane della forma qui
considerata: l'argomento della sommatoria nell'esponenziale in (\ref{eq:K_discr}) diventa
$$
p_i\dot q_i-\frac{p_i^2}{2m}-V(q_i)
$$
si completa il quadrato
$$
p_i\dot q_i-\frac{p_i^2}{2m}=-\frac{1}{2m}(p_i-m\dot q_i)^2 +\frac{m}{2}\dot q_i^2:=
-\frac{\bar p_i^2}{2m}+\frac{m}{2}\dot q_i^2
$$
e la parte dipendente dai $p_i$ di (\ref{eq:K_discr}) diventa
$$
I=\int\prod_{k=0}^{N-1}\frac{d^d\bar p_k}{(2\pi\hbar)^d}\exp\left\{\sum_{i=0}^{N-1}\frac{-i}{2m\hbar}
\bar p_i^2\Delta t\right\} = \left(\int_{-\infty}^{\infty}\frac{dp}{2\pi\hbar}\exp\left\{\frac{-i}{2m\hbar} p^2\Delta t\right\}\right)^{Nd}
$$
che è uno degli integrali calcolati per la particella libera e vale
$$I = \left(\frac{m}{2\pi i\hbar\Delta t}\right)^{Nd/2} = \left(\frac{1}{\mathcal{N}}\right)^{Nd}$$
Inserendo questo risultato nell'espressione di $K$ si trova
$$
K=\frac{1}{\mathcal{N}^d}\int\prod_{i=1}^{N-1}\frac{d^dq_i}{\mathcal{N}^d}\exp\left\{\frac{i}{\hbar}
\sum_{i=0}^{N-1}\Delta t\left(\frac{1}{2}m\dot q_i^2-V(q_i)\right)\right\}
$$
Passando al continuo e ignorando le normalizzazioni (che, come si vedrà, sarà facile fissare caso per
caso), si trova
\begin{equation}
\boxed{K(q',t',q,t)=\int\mathcal{D}q\exp\left\{\frac{i}{\hbar}\int_{t}^{t'}dtL(q(t),\dot q(t))
\right\}=\int\mathcal{D}q\exp\left\{\frac{i}{\hbar}S[q]\right\}}
\label{eq:K}
\end{equation}

\textbf{Oss.} Questa definizione di $K$ ne evidenzia esattamente il legame con l'azione e i cammini:
il valore di $K$ è l'integrale su tutti i possibili cammini che uniscono $q'$ e $q$, pesati con
\textit{il valore dell'azione sul cammino}.
\subsection{Equivalenza con l'equazione di Schrödinger}
